\chapter{Synthetic extension spectral sequences}
\label{chap:synthetic-extension-ss}

This chapter defines weightwise extension spectral sequences, the basic $\lambda$--$\rho$--$\delta$ triangle, and coherent solution towers.  Comparison theorems are deferred to the internal-results layer.

\section{Weightwise extension spectral sequences}

% SOURCE: aimpaper/main.tex:998-1007.
\begin{definition}[Synthetic extension spectral sequence]
  \label{def:synthetic-extension-ss}
  \uses{def:synthetic-adams-ss,def:f-extension-ss}
  \notready
  For a weight-preserving map $f:X\to Y$ of finite synthetic spectra equipped
  with convergence-and-detection witnesses for their synthetic Adams
  sequences, filter
  $0\to\pi_{*,*}X\to\pi_{*,*}Y\to0$ by Adams filtration.  The resulting
  $f$-ESS has
  \[
    {}^fE_0^{s,t,w}\cong E_\infty^{s,t,w}(X)\oplus
    E_\infty^{s,t,w}(Y),
    \qquad d_n^f:(s,t,w)\mapsto(s+n,t+n,w).
  \]
\end{definition}

% SOURCE: aimpaper/main.tex:999-1007, assertion that all Section 2 results remain valid.
\begin{proposition}[Weightwise transport of the ESS calculus]
  \label{prop:synthetic-ess-weightwise-transport}
  \uses{def:synthetic-extension-ss,prop:fextension-iff-detected,prop:fextension-inessential-iff-higher,prop:no-crossing-iff-uniform-detection,thm:fess-commutative-square,prop:fess-exact-middle-boundary}
  \notready
  Every definition and theorem of Chapter~\ref{chap:extension-ss} holds in a
  fixed synthetic weight, provided every map and comparison preserves that
  weight.  In particular, essentiality, crossings, the representative
  criterion, square propagation, and exact-middle boundary detection retain
  their statements with a third grading coordinate.
\end{proposition}

\begin{proof}
  Fix $w$ and regard the bigraded homotopy groups in that weight as filtered
  graded groups.  Reuse the two-term filtered-complex construction and each
  proof from Chapter~\ref{chap:extension-ss}; all maps preserve $w$, so no
  additional reindexing occurs.  Reassemble the family over all weights.
\end{proof}

% SOURCE: aimpaper/main.tex:1981-2081, representative equations used in the
% page-restriction argument; project packaging of their inverse system.
\begin{definition}[Finite ESS solution torsors and coherent towers]
  \label{def:synthetic-ess-solution-tower}
  \uses{def:synthetic-extension-ss,def:fess-cycles-boundaries,cor:fess-map-of-stable-pages,prop:lambda-quotient-restriction-coherence}
  \notready
  Let $F_m:A_m\to B_m$, $m\ge m_0$, be compatible maps of finite
  $\lambda$-quotients, with restriction maps from level $m+1$ to level $m$.
  Fix a synthetic weight, an extension length $n$, and compatible source and
  target classes $x_m,y_m$.  The \emph{solution fiber}
  $\operatorname{Sol}_m^n(x_m,y_m)$ consists of representative-level
  witnesses for
  \[
    d_n^{F_m}(x_m)=y_m:
  \]
  a filtered representative of $x_m$, a representative of the prescribed
  complete target coset of $y_m$, and a sum of shorter $F_m$-ESS images whose
  representative equation realizes that coset.  Two witnesses are identified
  exactly when their source representatives differ by a filtered boundary and
  their target equations differ by the same shorter-image correction.

  When nonempty, $\operatorname{Sol}_m^n(x_m,y_m)$ is a torsor under the
  abelian \emph{difference group} $K_m$ of changes preserving the source class
  and complete target coset.  Proposition~\ref{prop:lambda-quotient-restriction-coherence}
  makes the finite restriction maps an inverse system; stable-page naturality
  gives equivariant maps
  \[
    \operatorname{res}_{m+1,m}:\operatorname{Sol}_{m+1}^n
      \longrightarrow\operatorname{Sol}_m^n,
    \qquad K_{m+1}\longrightarrow K_m.
  \]
  A \emph{coherent solution tower} is a sequence
  $u_m\in\operatorname{Sol}_m^n$ satisfying
  $\operatorname{res}_{m+1,m}(u_{m+1})=u_m$ for every $m\ge m_0$.
\end{definition}

\begin{proof}
  Use the filtered two-term complex defining the synthetic ESS.  A choice of
  filtered source lift and target correction is precisely a solution of its
  representative equation.  Subtracting two solutions gives the stated
  difference group, and addition acts freely and transitively.  The compatible
  quotient squares induce the restriction maps and their equivariance.
\end{proof}

% SOURCE: aimpaper/main.tex:1009--1021 defines the quotient triangles;
% Burklund--Xu Construction 7.7, reference/BurklundXu/paper.txt:3022--3035,
% supplies the finite tower and its cofiber sequences.  Naturality and the
% identity/composition laws below are project-level consequences of the finite
% cofiber presentations, not additional external algebra input.
\begin{proposition}[Coherence of quotient restrictions]
  \label{prop:lambda-quotient-restriction-coherence}
  \uses{def:synthetic-lambda-quotient,def:cofiber-triangle,def:lambda-rho-delta-triangle}
  \notready
  For every map of classical spectra $f:X\to Y$ in the stable homotopy
  context (so $\nu f:\nu X\to\nu Y$ is the induced synthetic map), put
  $Q_j(X):=\nu X/\lambda^j$ for each finite $j\ge1$ and write
  $\bar f_j:Q_j(X)\to Q_j(Y)$ for the map induced by $\nu f$.  For finite
  $1\le i\le j$, the quotient restriction maps are natural:
  \[
    \rho^Y_{i,j}\circ\bar f_j
      =\bar f_i\circ\rho^X_{i,j}.
  \]
  They satisfy the identity and composition laws
  \[
    \rho^X_{i,i}=\operatorname{id},\qquad
    \rho^X_{k,i}\circ\rho^X_{i,j}=\rho^X_{k,j}
    \quad(1\le k\le i\le j).
  \]
  For every finite pair $0<i<j$, the maps
  $\Sigma^{0,-i}\bar f_{j-i}$, $\bar f_j$, $\bar f_i$, and
  $\Sigma^{1,-i}\bar f_{j-i}$ form a morphism from the finite
  $\lambda$--$\rho$--$\delta$ triangle for $X$ at $(i,j)$ to the
  corresponding triangle for $Y$.  Thus multiplication, restriction, and
  connecting maps commute in the finite tower, naturally in $f$.  This
  proposition asserts no map of $E_\infty$- or $\mathrm{CAlg}$-objects; the
  $m=\infty$ endpoint is handled separately by the complete-tower statements.
\end{proposition}

\begin{proof}
  Write $\mu^X_j:\Sigma^{0,-j}\nu X\to\nu X$ for multiplication by
  $\lambda^j$ (and put $\mu^X_0=\operatorname{id}_{\nu X}$).  For each finite
  $j$, the square
  \[
    \begin{array}{ccc}
      \Sigma^{0,-j}\nu X & \xrightarrow{\mu^X_j} & \nu X \\
      {\scriptstyle\Sigma^{0,-j}\nu f}\downarrow && \downarrow{\scriptstyle\nu f} \\
      \Sigma^{0,-j}\nu Y & \xrightarrow{\mu^Y_j} & \nu Y
    \end{array}
  \]
  commutes by naturality of multiplication by $\lambda^j$.  Applying the
  chosen cofiber construction gives $\bar f_j$.  For $i\le j$, define the shifted
  factor map
  \[
    q^X_{i,j}:=\Sigma^{0,-i}\mu^X_{j-i}:
      \Sigma^{0,-j}\nu X\longrightarrow\Sigma^{0,-i}\nu X,
    \qquad \mu^X_i\circ q^X_{i,j}=\mu^X_j.
  \]
  The square with horizontal maps $\mu^X_j,\mu^X_i$ and vertical maps
  $q^X_{i,j},\operatorname{id}_{\nu X}$ induces $\rho^X_{i,j}$ by the
  chosen cofiber construction.  The identities
  $q^X_{i,i}=\operatorname{id}$ and
  $q^X_{k,i}\circ q^X_{i,j}=q^X_{k,j}$ induce the displayed identity and
  composition laws.  Naturality of the $q^X_{i,j}$ squares with respect to
  $\nu f$, followed by the chosen triangle equivalences, gives the morphism
  of finite distinguished triangles with the four components displayed above.
  This is a finite cofiber calculation and uses no multiplicative algebra
  structure.
\end{proof}

\section{The elementary lambda and rho extensions}

% SOURCE: aimpaper/main.tex:1009-1021, Notation not:deltaandrho.
\begin{definition}[$\lambda$--$\rho$--$\delta$ quotient triangle]
  \label{def:lambda-rho-delta-triangle}
  \uses{def:synthetic-lambda-quotient,def:cofiber-triangle}
  \notready
  For $0<n<m\le\infty$, define the distinguished triangle
  \[
    \Sigma^{0,-n}\nu X/\lambda^{m-n}
      \xrightarrow{\lambda^n}\nu X/\lambda^m
      \xrightarrow{\rho_{n,m}}\nu X/\lambda^n
      \xrightarrow{\delta_{n,m}}
      \Sigma^{1,-n}\nu X/\lambda^{m-n}.
  \]
  For $m=\infty$, the first two terms are
  $\Sigma^{0,-n}\nu X$ and $\nu X$.  Translation in the synthetic category is
  suspension by $S^{1,0}$.
\end{definition}
